
\section{Canvasflow Guarantees}\label{sec:canvasflow-guarantees}

\textit{Normative language per \href{https://tools.ietf.org/html/rfc2119}{RFC~2119}.}

\subsection{Purpose}

This section defines the commitments Canvasflow makes to every tenant who has completed
onboarding and whose \texttt{status} is \texttt{active}. These guarantees define
Canvasflow's obligations under the integration contract and bound the expectations a
tenant may reasonably hold of the Canvasflow Resource Server.

\subsection{Signature Validation on Every Request}

Canvasflow WILL validate the cryptographic signature and all required claims of every
access token received at the GraphQL API on every request, without exception. No request
is granted content access without passing the full validation sequence.
 There are no bypass paths, administrative overrides, or ``trusted subnet'' 
 exceptions to token validation.

\subsection{Content Gating}

Canvasflow WILL gate GraphQL content strictly by the verified entitlement claim --- and
nothing else. Specifically:



\begin{itemize}
  \item Content Canvasflow classifies as \textbf{free} WILL be returned to any bearer of
    a valid, verified access token, regardless of the effective entitlement set (even if
    the entitlement set is empty).

  \item Content Canvasflow classifies as \textbf{premium} WILL be returned only when the
    matching \texttt{cf:<type>:<id>} value is present in the effective entitlement set
    (computed per \hyperref[sec:mandatory-value-format]{Entitlement Mapping \S4.2}).

  \item Canvasflow WILL NOT invent, infer, or apply any additional authorization logic
    beyond what the signed token encodes. If a resource identifier appears in the
    verified entitlement set, access is granted. If it does not appear, access is not
    granted. No exceptions.
\end{itemize}

\subsection{JWKS Caching Behavior}

Canvasflow's JWKS caching behavior for JWKS-path tenants:

\begin{itemize}
  \item JWKS documents are fetched and cached \textbf{server-side only} with a cache TTL
    of one hour. JWKS key material is never exposed to clients or included in API
    responses.

  \item An incoming token whose \texttt{kid} is not present in the current cache WILL
    trigger an \textbf{immediate re-fetch} of the tenant's JWKS endpoint. The token is
    validated against the freshly fetched key set before being accepted or rejected.

  \item If the JWKS re-fetch fails (network error, timeout, or non-200 response),
    Canvasflow WILL NOT fall back to the stale cached key set. The token WILL be rejected
    with 401. This policy prevents a degraded JWKS endpoint from becoming a persistent
    authentication bypass.

  \item Canvasflow WILL NOT cache JWKS data in a way that persists beyond the one-hour
    TTL unless a re-fetch confirms the key is still valid.
\end{itemize}

\subsection{HS256 Secret Handling}

For HS256-path tenants, Canvasflow commits to the following:

\begin{itemize}
  \item Shared secrets are stored \textbf{encrypted at rest} using AES-256-GCM at the
    application layer, with a KMS-managed encryption key. Database-level encryption alone
    is not the sole protection.

  \item Secrets WILL NEVER appear in server logs, error responses, debugging output, or
    any client-visible channel, under any circumstances.

  \item Canvasflow WILL notify the affected tenant within \textbf{24~hours} of any
    confirmed or suspected compromise of a stored HS256 secret, including disclosure of
    the nature and scope of the compromise and the steps taken to mitigate it.
\end{itemize}

\subsection{Generic Error Responses}

All token validation failures WILL return the same generic 401 response:

\begin{lstlisting}
HTTP 401 Unauthorized
WWW-Authenticate: Bearer error="invalid_token"
\end{lstlisting}

The response body WILL contain only a generic error message. It WILL NOT reveal:
\begin{itemize}
  \item Which validation step failed (e.g., ``expired token'' vs.\ ``invalid signature'')
  \item Which tenant's token was presented
  \item Which claim was absent or malformed
  \item Which signing path was attempted
\end{itemize}

Detailed failure reasons are logged server-side and correlated by \texttt{jti} (when
extractable). These logs are accessible to Canvasflow support for debugging with the
tenant's consent; they are never included in client-facing responses.

This policy is a security requirement. Specific validation failure messages enable
targeted token-forgery attacks by providing feedback about which properties of a crafted
token are correct.

\subsection{Resource Identifier Stability}

Canvasflow resource identifiers (e.g., \texttt{cf:issue:123}) assigned to a tenant
during Phase~1 of onboarding WILL NOT change after go-live without:

\begin{itemize}
  \item Prior written notice to the tenant's designated technical lead
  \item A mutually agreed migration window sufficient for the tenant to update their IdP
    hook mapping table
\end{itemize}

A change to resource identifiers without advance notice would silently break the tenant's
entitlement mapping and deprive entitled users of content. Canvasflow treats resource
ID stability as a hard commitment.

\subsection{No Runtime Calls to the Tenant IdP}

Canvasflow WILL NOT call the tenant's IdP at API request time for any purpose other than
JWKS cache refresh (triggered by an unknown \texttt{kid}). The tenant's IdP availability
is NOT a dependency for serving already-issued access tokens. A tenant IdP outage WILL
NOT cause Canvasflow to reject valid, unexpired tokens that were issued before the outage
began.

Canvasflow will not perform token introspection, revocation checks, or any other
real-time verification against the tenant's Authorization Server on the request path.

\subsection{No Token Issuance by Canvasflow}

Canvasflow WILL NEVER issue, mint, modify, extend, or re-sign tokens on a tenant's
behalf. All token issuance is the exclusive responsibility of the tenant's Authorization
Server. Canvasflow's cryptographic role is limited to verification.
